
\section{Method}
\label{sec:method}

\textbf{ACR-VLA} improves the actions of a frozen flow-matching VLA at inference time through a sample-then-correct pipeline (Figure~\ref{fig:pipeline}): draw $K$ action samples, select the best candidate via conformalized scoring, then apply a learned residual correction bounded in magnitude. The base VLA is never retrained.

\subsection{Problem Setup}
\label{sec:problem_setup}

Let $\pi_\theta$ be a frozen flow-matching VLA that, given an observation $o_t$ (image + language instruction + proprioceptive state), generates action chunks of length $C{=}50$ in $\mathbb{R}^{d \times C}$ with $d{=}6$ (3 translation + 3 rotation). Under receding-horizon execution, only the first $H{=}5$ steps of each chunk influence the robot before re-planning, so we discard the remainder. The gripper dimension is handled separately via direct thresholding and is excluded from the conformal pipeline.

At each control step, $\pi_\theta$ can be queried multiple times via independent flow-matching rollouts~\cite{lipman2023flow}, yielding $K$ candidate action chunks. Our goal is to exploit these $K$ samples to produce a single action that is more accurate than any individual sample or their aggregate statistics, without modifying $\pi_\theta$.

\subsection{Multi-Sample Generation and Conformal Selection}
\label{sec:selection}

\paragraph{Sampling and truncation.}
We draw $K{=}10$ action chunks from $\pi_\theta$ via independent flow-matching rollouts. Each chunk is truncated to the first $H{=}5$ steps, yielding samples $\{a_h^{(k)}\}_{k=1}^K$ with $a_h^{(k)} \in \mathbb{R}^d$ for each horizon step $h = 1, \dots, H$.

\paragraph{Shrinkage covariance estimation.}
For each horizon step $h$, we compute the sample mean $\boldsymbol{\mu}_h = \frac{1}{K}\sum_{k=1}^K a_h^{(k)}$ and the Ledoit--Wolf shrinkage covariance~\cite{ledoit2004well}:
\begin{equation}
  \boldsymbol{\Sigma}_h = (1 - \alpha_{\mathrm{LW}})\,\hat{\boldsymbol{\Sigma}}_{\mathrm{emp}} + \alpha_{\mathrm{LW}}\,\frac{\mathrm{tr}(\hat{\boldsymbol{\Sigma}}_{\mathrm{emp}})}{d}\,\mathbf{I},
  \label{eq:shrinkage}
\end{equation}
where $\hat{\boldsymbol{\Sigma}}_{\mathrm{emp}}$ is the empirical covariance and $\alpha_{\mathrm{LW}}$ is the analytically optimal shrinkage intensity. With dimension-to-sample ratio $d/K = 0.6$, the eigenvalues of the raw sample covariance suffer 30--40\% estimation error; shrinkage reduces this to below 5\%~\cite{ledoit2004well}.

\paragraph{Mahalanobis nonconformity score.}
For each sample $a_h^{(k)}$ at horizon step $h$, we compute the Mahalanobis nonconformity score:
\begin{equation}
  s_h(a) = (a - \boldsymbol{\mu}_h)^\top \boldsymbol{\Sigma}_h^{-1} (a - \boldsymbol{\mu}_h).
  \label{eq:mahal}
\end{equation}
This score captures cross-DOF covariance structure: a sample that deviates along the principal axis of variation receives a lower score than one that deviates along a tightly concentrated direction~\cite{braun2025cgs}.

\paragraph{Argmin selection.}
We select the sample with the lowest total nonconformity across all horizon steps:
\begin{equation}
  a_{\mathrm{base}} = \argmin_{k \in \{1,\dots,K\}} \sum_{h=1}^{H} s_h\bigl(a_h^{(k)}\bigr).
  \label{eq:argmin}
\end{equation}
This selection procedure uses Mahalanobis-based scoring~\cite{braun2025cgs} to identify the sample closest to the consensus under the empirical covariance structure. The selected $a_{\mathrm{base}} \in \mathbb{R}^{d \times H}$ is the starting point for residual correction.

\subsection{Conformal Residual Module}
\label{sec:crm}

The Conformal Residual Module (CRM) is a lightweight MLP that refines the selected action $a_{\mathrm{base}}$ by predicting a small residual correction conditioned on the current observation.

\paragraph{Architecture.}
CRM is a 2-layer MLP with hidden dimension 256 (${\sim}$N parameters).% TODO: exact param count pending from checkpoint
Its input is the frozen observation embedding $\mathbf{z}_{\mathrm{obs}}$ extracted from the base VLA's vision-language encoder. The output is a residual $\boldsymbol{\delta} \in \mathbb{R}^{d \times H} = \mathbb{R}^{30}$, and the corrected action is:
\begin{equation}
  \mathbf{a}^* = a_{\mathrm{base}} + \mathrm{clip}\bigl(\boldsymbol{\delta},\; \|\boldsymbol{\delta}\|_2 \leq \varepsilon\bigr),
  \label{eq:crm}
\end{equation}
where $\varepsilon$ is a norm bound hyperparameter (Section~\ref{sec:magnitude}).

\paragraph{Observation-only conditioning.}
CRM receives only the observation embedding $\mathbf{z}_{\mathrm{obs}}$ and does \emph{not} receive per-DOF conformal statistics, uncertainty features, or any CGS-derived quantities as input. This is a deliberate design choice validated by ablation (Section~\ref{sec:ablation}): observation features already encode sufficient task-relevant information---object identity, spatial layout, gripper-object relationship---for predicting corrective residuals. Per-DOF conformal statistics, while intuitively richer, add noise at the $K{=}10$ sample budget where covariance estimation is inherently noisy ($d/K = 0.6$). In Section~\ref{sec:ablation}, we compare against CRM variants conditioned on per-DOF features and scalar conformal scores, and find that observation-only conditioning is uniquely effective.

\paragraph{Training.}
CRM is trained offline with an $L_2$ loss on demonstration trajectories from LIBERO~\cite{liu2023libero}:
\begin{equation}
  \mathcal{L} = \bigl\|\mathbf{a}^* - \mathbf{a}_{\mathrm{expert}}\bigr\|_2^2.
  \label{eq:loss}
\end{equation}
The training data comprises 26{,}737 precomputed samples derived from 10 demonstrations per task across 20 LIBERO tasks. For each demonstration step, we run $K{=}10$ forward passes through the frozen SmolVLA~\cite{shukor2025smolvla}, compute the CGS scores, and apply argmin selection (\ref{eq:argmin}) to obtain $a_{\mathrm{base}}$. The expert action $\mathbf{a}_{\mathrm{expert}}$ is the corresponding demonstration action.

We apply two forms of data augmentation during training: (i)~additive Gaussian noise on the mean-action vector ($\sigma \in [0.02, 0.1]$) and (ii)~noise on the uncertainty features ($\sigma \approx 0.006$). Both augmentations improve generalization to the distribution of actions encountered during closed-loop evaluation, where compounding errors cause the state distribution to diverge from the demonstration distribution~\cite{ross2011dagger}.

\paragraph{Base VLA.}
We instantiate ACR-VLA on SmolVLA~\cite{shukor2025smolvla}, a 450M-parameter flow-matching VLA. SmolVLA generates actions via learned flow matching~\cite{lipman2023flow}, producing stochastic samples that naturally support multi-sample inference. The method is architecture-agnostic and applies to any flow-matching VLA that supports repeated sampling.

\subsection{Magnitude Control}
\label{sec:magnitude}

The norm bound $\varepsilon$ in~(\ref{eq:crm}) is critical: it determines how much the CRM can deviate from the conformal-selected action.

\paragraph{Unbounded correction fails catastrophically.}
Without clipping, the natural CRM output has $\|\boldsymbol{\delta}\|_2 \approx 1.42$. At this magnitude, the corrected action deviates so far from the base VLA's output that compounding error across the trajectory destroys task performance entirely---yielding 0\% success on LIBERO-Object. This failure mode is analogous to the compounding error problem in imitation learning~\cite{ross2011dagger}, but manifests at inference time: each over-corrected action shifts the state distribution further from the regime where the CRM was trained, amplifying subsequent correction errors.

\paragraph{Small corrections are effective.}
Constraining $\varepsilon$ to the range $[0.05, 0.2]$ yields stable, beneficial corrections. We use $\varepsilon = 0.1$ as the default. At this magnitude, the CRM applies corrections small enough to remain within the basin of attraction of the base policy's behavior, while large enough to shift actions toward the expert demonstrations.

\paragraph{Connection to residual policy learning.}
The necessity of magnitude control connects ACR-VLA to the broader residual policy literature. Silver et al.~\cite{silver2018residual} introduced residual policy learning with zero-initialization so the learned correction starts as a no-op. TRANSIC~\cite{jiang2024transic} uses a learned binary gate conditioned on current state to decide whether to apply a sim-to-real residual. PLD~\cite{xiao2026pld} bounds residual actor corrections within a fixed range during its probe-learn-distill pipeline. ACR-VLA's approach---fixed $L_2$ norm clipping of the CRM output---is the simplest instantiation: no learned gating, no scheduling, just a hard constraint that prevents the correction from overwhelming the base policy.

\paragraph{Differential sensitivity across CRM variants.}
Observation-conditioned CRM (obs\_only) maintains 35--36\% success on LIBERO-Spatial across the full range $\varepsilon \in [0.05, 0.3]$, while per-DOF-conditioned CRM collapses from 35\% at $\varepsilon = 0.1$ to 26\% at $\varepsilon = 0.3$. This robustness gap provides further evidence that observation features yield more stable conditioning than per-DOF conformal statistics at the $K{=}10$ sample budget (Section~\ref{sec:ablation}).

\subsection{Inference Protocol}
\label{sec:inference_protocol}

Each control step proceeds as follows:
\begin{enumerate}[nosep]
  \item Draw $K{=}10$ action chunks from the frozen VLA $\pi_\theta$ and truncate to $H{=}5$ steps.
  \item Compute shrinkage covariance $\boldsymbol{\Sigma}_h$ and Mahalanobis scores $s_h$ for each sample.
  \item Select $a_{\mathrm{base}}$ via argmin over total nonconformity~(\ref{eq:argmin}).
  \item Compute residual $\boldsymbol{\delta} = \mathrm{CRM}(\mathbf{z}_{\mathrm{obs}})$ and clip to $\|\boldsymbol{\delta}\|_2 \leq \varepsilon$.
  \item Execute the first action of $\mathbf{a}^* = a_{\mathrm{base}} + \mathrm{clip}(\boldsymbol{\delta})$.
  \item Observe new state and re-plan from step~1.
\end{enumerate}
The computational overhead consists of $K{=}10$ extra forward passes through the base VLA (shared with any multi-sample selection method such as ReconVLA~\cite{chen2026reconvla}) plus ${\sim}$3\,ms for CRM inference. The base VLA is never modified.
